\section{Dimensión, profundidad y regularidad en los preesquemas localmente
noetherianos}

Este parágrafo se limita a retomar, en lenguaje geométrico y con diversos
complementos de naturaleza técnica, las nociones y los resultados de álgebra
conmutativa expuestos en el capítulo \textbf{0}, \S\S~16 y 17.

\subsection{Dimensión de los preesquemas.}

\begin{env}[5.1.1]
\label{IV.2.5.1.1-fr}
Sean $A$ un anillo y $\mathfrak J$ un ideal de $A$. Recordemos
$(\mathbf{0},16.1)$ que las definiciones de la teoría de la dimensión de un anillo
$(\mathbf{0},16.1.1)$ y las de la teoría de la dimensión combinatoria de los espacios
topológicos $(\mathbf{0},14.1.2)$ dan las relaciones siguientes:
\begin{align}
\dim(\operatorname{Spec}(A))&=\dim(A)
\tag{5.1.1.1}\label{IV.2.5.1.1.1-fr}\\
\dim(V(\mathfrak J))&=\dim(A/\mathfrak J)
\tag{5.1.1.2}\label{IV.2.5.1.1.2-fr}\\
\operatorname{codim}(V(\mathfrak J),\operatorname{Spec}(A))
&=\operatorname{ht}(\mathfrak J)
\tag{5.1.1.3}\label{IV.2.5.1.1.3-fr}
\end{align}
\par\noindent\textbf{(5.1.1.4)} $\operatorname{Spec}(A)$ es un espacio
catenario $\Longleftrightarrow A$ es un anillo catenario.
\par\noindent\textbf{(5.1.1.5)} $\operatorname{Spec}(A)$ es equidimensional
$\Longleftrightarrow A$ es equidimensional
$\Longleftrightarrow$ los ideales primos minimales $\mathfrak p_\alpha$ de
$A$ son tales que las dimensiones de los anillos $A/\mathfrak p_\alpha$ son
todas iguales.
\par\noindent\textbf{(5.1.1.6)} $\operatorname{Spec}(A)$ es
equicodimensional $\Longleftrightarrow A$ es equicodimensional
$\Longleftrightarrow$ todos los ideales maximales de $A$ tienen la misma
altura. Se recuerda que un anillo noetheriano $A$ se llama
\emph{biequidimensional} si $\operatorname{Spec}(A)$ es biequidimensional, es
decir, si $\operatorname{Spec}(A)$ es noetheriano y $A$ es equidimensional,
equicodimensional, catenario y de dimensión finita.
\end{env}

